\documentclass[../../main.tex]{subfiles}
\begin{document}

{\bf[解]}

\paragraph{(1)}

% TODO: 双曲線，接戦とtriangle OQR
\begin{figure}[htb]
  \centering
  \begin{tikzpicture}[scale=1.7, >=stealth, every node/.style={font=\small}]

    % --- 1. 定数・幾何パラメータ ---
    % a = 1.6, b = 1.2
    % 接点 P(X, Y): X = a*cosh(t), Y = b*sinh(t) (t = 0.55 とする)
    \def\aVal{1.6}
    \def\bVal{1.2}
    \def\tVal{0.55}

    \pgfmathsetmacro{\xP}{\aVal * cosh(\tVal)} % ≈ 1.849
    \pgfmathsetmacro{\yP}{\bVal * sinh(\tVal)} % ≈ 0.694

    % A = X/a + Y/b = cosh(t) + sinh(t) = e^t
    % B = X/a - Y/b = cosh(t) - sinh(t) = e^(-t)
    \pgfmathsetmacro{\valA}{exp(\tVal)}
    \pgfmathsetmacro{\valB}{exp(-\tVal)}

    % 交点 Q(a/A, -b/A) : y = -b/a x との交点
    % 交点 R(a/B, b/B)  : y =  b/a x との交点
    \coordinate (O) at (0, 0);
    \coordinate (P) at (\xP, \yP);
    \coordinate (Q) at ({\aVal/\valA}, {-\bVal/\valA});
    \coordinate (R) at ({\aVal/\valB}, {\bVal/\valB});

    % --- 2. 三角形 OQR の塗りつぶし ---
    \fill[pattern=north east lines, pattern color=blue!65, opacity=0.7]
    (O) -- (Q) -- (R) -- cycle;
    \node[draw=blue!85!black, rounded corners=2pt, inner sep=2.5pt, font=\footnotesize\bfseries, blue!90!black]
    at (1.1, 0.45) {$\triangle OQR$};

    % --- 3. 漸近線 y = ± (b/a) x ---
    \draw[dashed, gray!60, thick] (-2.2, {-2.2*\bVal/\aVal}) -- (3.2, {3.2*\bVal/\aVal})
    node[above right=-1pt, font=\scriptsize, gray!90!black] {$y = \frac{b}{a}x$};
    \draw[dashed, gray!60, thick] (-2.2, {2.2*\bVal/\aVal}) -- (3.2, {-3.2*\bVal/\aVal})
    node[below right=-1pt, font=\scriptsize, gray!90!black] {$y = -\frac{b}{a}x$};

    % --- 4. 座標軸 ---
    \draw[->, thick, gray!80!black] (-2.4, 0) -- (3.4, 0) node[right] {$x$};
    \draw[->, thick, gray!80!black] (0, -2.0) -- (0, 2.2) node[above] {$y$};
    \node[below left=2pt] at (O) {$O$};

    % --- 5. 双曲線 x^2/a^2 - y^2/b^2 = 1 ---
    % 右側の分枝
    \draw[very thick, blue!85!black, domain=-1.2:1.2, samples=60]
    plot ({\aVal*cosh(\x)}, {\bVal*sinh(\x)});
    % 左側の分枝（参考）
    \draw[thick, gray!45, domain=-1.2:1.2, samples=60]
    plot ({-\aVal*cosh(\x)}, {\bVal*sinh(\x)});
    \node[blue!90!black, font=\footnotesize, right] at (2.4, 1.6) {$\frac{x^2}{a^2} - \frac{y^2}{b^2} = 1$};

    % --- 6. 点 P における接線 l (Q と R を結ぶ直線) ---
    % 接線分およびその延長
    \draw[line width=1.6pt, red!85!black] ($ (Q)!-0.2!(R) $) -- ($ (Q)!1.2!(R) $);
    \node[red!90!black, font=\footnotesize, above right] at ($ (Q)!1.15!(R) $) {接線 $l$};

    % 接点 P の座標投影破線
    \draw[dotted, gray!60] (\xP, 0) -- (P) -- (0, \yP);
    \node[below=2pt, font=\scriptsize] at (\xP, 0) {$X$};
    \node[left=2pt, font=\scriptsize]  at (0, \yP) {$Y$};

    % 頂点 (a, 0)
    \node[below left=1pt, font=\scriptsize] at (\aVal, 0) {$a$};
    \fill (\aVal, 0) circle (0.8pt);

    % --- 7. 主要点のプロットとラベル ---
    \fill (O) circle (0.8pt);
    \fill[red!85!black]  (P) circle (1.2pt) node[above left=1pt, red!95!black] {$P(X, Y)$};
    \fill[teal!85!black] (Q) circle (1.2pt) node[below right=1pt, teal!95!black] {$Q$};
    \fill[teal!85!black] (R) circle (1.2pt) node[above left=1pt, teal!95!black] {$R$};

  \end{tikzpicture}
  \caption{双曲線 $\dfrac{x^2}{a^2} - \dfrac{y^2}{b^2} = 1$ の接線 $l$ と漸近線の交点 $Q, R$ がなす三角形 $\triangle OQR$}
  \label{fig:hyperbola_tangent_triangle}
\end{figure}

$2$本の漸近線は
\begin{align}
  y=\pm\dfrac{b}{a}x \label{eq:1}
\end{align}
であり，$P(X,Y)$での双曲線の接線$l$は
\begin{align}
  \dfrac{Xx}{a^2}-\dfrac{Yy}{b^2}=1 \label{eq:2}
\end{align}
である．\cref{eq:1,eq:2}の交点はまず$y$を消去して
\begin{align}
  \dfrac{Xx}{a^2}\mp\dfrac{Y}{b^2}\dfrac{b}{a}x =1 \\
  \dfrac{Xx}{a^2}\mp\dfrac{Yx}{ab} =1 \\
  x = \dfrac{a}{X/a \mp Y/b}
\end{align}
だから対応して$y$は
\begin{align}
  y = \pm\dfrac{b}{a} \dfrac{a}{X/a \mp Y/b} = \dfrac{\pm b}{X/a \mp Y/b}
\end{align}
である．ここで簡単のために
\begin{align}
  A=\dfrac{X}{a}+\dfrac{Y}{b} , B=\dfrac{X}{a}-\dfrac{Y}{b} \label{eq:3}
\end{align}
とおくと，これらは
\begin{align}
  Q\left(\dfrac{a}{A},\dfrac{-b}{A}\right),R\left(\dfrac{a}{B},\dfrac{b}{B}\right)
\end{align}
である．したがってサラスの公式から三角形$OQR$の面積は
\begin{align}
  S&=\dfrac{1}{2}\left|\dfrac{ab}{AB}+\dfrac{ab}{AB}\right|  \nonumber\\
  &=\dfrac{ab}{AB}  \nonumber\\
  &=ab \ \ \ \left(\because AB=\dfrac{X^2}{a^2}-\dfrac{Y^2}{b^2}=1\right)  \label{eq:4}
\end{align}
となる．これは$X$，$Y$によらない定数であるから題意は示された．$\Box$

\paragraph{(2)}

$p=e^t$とすると$t$が任意実数だから$p>0$である．\cref{eq:4}に代入して
\begin{align}
  S&=(5p^2+p^{-1})(p^2+p^{-1})\nonumber \\
  &=5p^4+6p+p^{-2} \label{eq:5}
\end{align}
であるから以下この最小値を求める．
一階微分を計算すると
\begin{align}
  \dfrac{dS}{dp}&=20p^3+6-\dfrac{2}{p^3} \\
  &=\dfrac{2}{p^3}(5p^3-1)(2p^3+1)
\end{align}
だから$S$の増減表は以下のようになる．

\begin{table}[htb]
  \centering
  \caption{$S(p)$ の増減表}
  \label{tab:sp_monotonicity}
  \renewcommand{\arraystretch}{1.5}
  \begin{tabular}{c|ccccc}
    $p$ & $(0)$ & $\dots$ & $\left(\dfrac{1}{5}\right)^{1/3}$ & $\dots$ \\ \hline
    $S'$ & & $-$ & $0$ & $+$ \\ \hline
    $S$  & & $\searrow$ & \text{極小} & $\nearrow$
  \end{tabular}
\end{table}

したがって$S$が最小になるのは$p=\left(\dfrac{1}{5}\right)^{1/3}$の時で
\cref{eq:5}に代入して求める最小値は
\begin{align}
  \min S
  &=S\left(p=\left(\dfrac{1}{5}\right)^{1/3}\right) \\
  &=5\left(\dfrac{1}{5}\right)^{4/3}+6\left(\dfrac{1}{5}\right)^{1/3}+5^{2/3} \\
  &=7\left(\dfrac{1}{5}\right)^{1/3}+5^{2/3}\cdots\text{(答)}
\end{align}
となる．

\newpage
\end{document}
